\Needspace{9\baselineskip}\begin{center}Über die durchschnittliche Anzahl der Geschlechter.\end{center}


\Needspace{7\baselineskip}\Article{301}
Wir kehren zu den binären Formen zurück, von denen noch mehrere bemerkenswerte Eigenschaften zu besprechen sind. Zunächst fügen wir einige Bemerkungen zur Anzahl der Geschlechter und Klassen in der eigentlich primitiven Ordnung hinzu, bei negativer Determinante in der positiven. Auf diese Ordnung beschränken wir der Kürze wegen die Untersuchung.

Die \textit{Anzahl der Geschlechter}, in die sich alle eigentlich primitiven Formen, bei negativer Determinante die positiven, einer gegebenen positiven oder negativen Determinante \(\pm D\) verteilen, ist immer \(1,2,4\) oder eine höhere Potenz von \(2\). Ihr Exponent hängt von den Faktoren von \(D\) ab und lässt sich aus den vorigen Untersuchungen vollständig im Voraus bestimmen. Da in der natürlichen Zahlenreihe Primzahlen mit mehr oder weniger zusammengesetzten Zahlen vermischt sind, nimmt die Anzahl der Geschlechter bei mehreren
\SourceLeaf{368}
aufeinanderfolgenden Determinanten \(\pm D,\pm(D+1),\pm(D+2)\) usw.\ bald zu, bald ab, und diese unregelmäßige Reihe scheint keine Ordnung zu haben. Addiert man aber die Anzahlen der Geschlechter für viele aufeinanderfolgende Determinanten
\[\pm D,\quad\pm(D+1)\;\ldots\;\pm(D+m)\]
und teilt die Summe durch die Anzahl der Determinanten, so erhält man die \textit{durchschnittliche Anzahl der Geschlechter}. Diese kann man etwa der mittleren Determinante \(\pm(D+\frac12m)\) zuordnen; sie bildet eine sehr regelmäßige Folge. Dabei setzen wir nicht nur voraus, dass \(m\) hinreichend groß ist, sondern auch, dass \(D\) noch viel größer ist, sodass das Verhältnis der äußersten Determinanten \(D,D+m\) nicht zu stark von der Gleichheit abweicht. Die Regelmäßigkeit dieser Folge ist so zu verstehen: Ist \(D^\prime\) viel größer als \(D\), so ist die durchschnittliche Anzahl der Geschlechter um die Determinante \(\pm D^\prime\) merklich größer als um \(D\). Unterscheiden sich dagegen \(D,D^\prime\) nicht zu stark, so sind auch die durchschnittlichen Anzahlen der Geschlechter um \(D\) und \(D^\prime\) nahezu gleich. Ferner findet man die durchschnittliche Anzahl der Geschlechter um die positive Determinante \(+D\) stets nahezu gleich der durchschnittlichen Anzahl um die negative, und zwar umso genauer, je größer \(D\) ist. Für einen kleinen Wert ist die erste etwas größer als die zweite. Das verdeutlichen die folgenden Beispiele aus einer Klassifikationstafel der binären Formen für mehr als \(4000\) Determinanten. Unter den hundert Determinanten von \(801\) bis \(900\) gibt es \(7\), denen nur ein Geschlecht entspricht, sowie \(32,52,8,1\), denen der Reihe nach \(2,4,8.\ 16\) Geschlechter entsprechen. Insgesamt ergeben sich also \(359\) Geschlechter und damit die durchschnittliche Anzahl \(=3{,}59\). Die hundert negativen Determinanten von \(-801\) bis \(900\) ergeben \(360\) Geschlechter. Alle folgenden Beispiele stammen von negativen Determinanten. Im \(16\). Hunderter, von \(-1501\) bis \(-1600\), findet man die durchschnittliche Anzahl \(3{,}89\); im \(25\). Hunderter \(4{,}03\); im \(51\). Hunderter \(4{,}24\); aus den sechshundert Determinanten \(-9401\;\ldots -10000\) ergibt sich \(4{,}59\). Diese Beispiele zeigen, dass die durchschnittliche Anzahl der Geschlechter viel langsamer wächst als die Determinanten selbst. Nach welchem Gesetz verläuft aber diese Folge? — Eine ziemlich schwierige theoretische Untersuchung, deren Darstellung hier zu ausführlich wäre, hat ergeben, dass die durchschnittliche Anzahl der Geschlechter um die Determinante \(+D\) oder \(-D\) sehr nahe durch die Formel
\[\alpha\log D+\beta\]

\SourceLeaf{369}
\begingroup\rightskip=0pt plus 3em\relax
gegeben wird. Dabei sind \(\alpha,\beta\) Konstanten, und zwar
\[\alpha=\frac{4}{\pi\pi}=0{,}4052847346\]
wobei \(\pi\) den halben Umfang eines Kreises mit dem Radius \(1\) bezeichnet, und
\[\beta=2\alpha g+3\alpha\alpha h-\frac16\alpha\log2=0{,}8830460462\]
wobei \(g\) die Summe der Reihe
\[\begin{aligned}&1-\log(1+1)+\frac12-\log(1+\frac12)\\&\quad+\frac13-\log(1+\frac13)+\text{etc.}=0{,}5772156649\end{aligned}\]
ist (siehe Euler, \textit{Inst. Calc. Diff.}, S.\ 444), während \(h\) die Summe der Reihe
\[\frac14\log2+\frac19\log3+\frac1{16}\log4+\text{etc.}\]
ist, die durch Näherung \(=0{,}9375482543\) gefunden wurde. Aus dieser Formel folgt, dass die durchschnittliche Anzahl der Geschlechter in einer arithmetischen Folge wächst, wenn die Determinanten in einer geometrischen Folge wachsen. Für \(D=850\frac12\), \(1550\frac12,2450\frac12,5050\frac12,9700\frac12\) ergeben sich die Werte \(3{,}617;3{,}86;4{,}046;4{,}339;4{,}604\), die von den oben angegebenen Durchschnittswerten wenig abweichen. Je größer die mittlere Determinante ist und je mehr Determinanten zur Berechnung des Durchschnitts dienen, desto weniger weicht er vom Wert der Formel ab. Mit dieser Formel lässt sich auch die Summe der Anzahlen der Geschlechter für aufeinanderfolgende Determinanten \(\pm D,\pm(D+1)\ldots\pm(D+m)\) sehr nahe bestimmen, wie verschieden die äußersten \(D,D+m\) auch sein mögen: Man berechnet die jeweils zugehörigen Durchschnittswerte und addiert sie. Diese Summe ist
\[=\alpha(\log D+\log(D+1)+\text{etc.}+\log(D+m))+\beta(m+1)\]
oder hinreichend genau
\[\begin{aligned}={}&\alpha((D+m)\log(D+m)-(D-1)\log(D-1))\\&+(\beta-\alpha)(m+1)\end{aligned}\]
So ergibt sich für die Summe der Geschlechterzahlen der Determinanten \(-1\) bis \(-100\) der Wert \(=234{,}4\), während sie tatsächlich \(233\) beträgt; ebenso von \(-1\) bis \(-2000\) der Wert \(=7116{,}6\), während sie \(7112\) beträgt; von \(-9001\) bis \(-10000\), wo sie \(4595\) beträgt, gibt die Formel \(4594{,}9\). Eine solche Übereinstimmung hätte man kaum erwarten können.
\par\endgroup
\SourceLeaf{370}

\Needspace{9\baselineskip}\begin{center}Über die durchschnittliche Anzahl der Klassen.\end{center}


\Needspace{7\baselineskip}\Article{302}
Hinsichtlich der \textit{Anzahl der Klassen}, wobei stets eigentlich primitive Klassen gemeint sind, bei negativer Determinante die positiven, verhalten sich positive Determinanten ganz anders als negative. Deshalb betrachten wir beide getrennt. Sie stimmen darin überein, dass bei gegebener Determinante alle Geschlechter gleich viele Klassen enthalten. Somit ist die Gesamtzahl der Klassen gleich dem Produkt aus der Anzahl der Geschlechter und der Anzahl der Klassen in jedem Geschlecht.

Zunächst zu den negativen Determinanten: Die Anzahl der Klassen für mehrere aufeinanderfolgende Determinanten \(-D,-(D+1),-(D+2)\) usw.\ bildet eine ebenso unregelmäßige Folge wie die Anzahl der Geschlechter. Die durchschnittliche Anzahl der Klassen, die keiner Definition bedarf, wächst dagegen sehr regelmäßig, wie die folgenden Beispiele zeigen. Die hundert Determinanten von \(-500\) bis \(-600\) ergeben \(1729\) Klassen, also die durchschnittliche Anzahl \(=17{,}29\). Ebenso findet man im \(15\). Hunderter die durchschnittliche Anzahl \(28{,}26\); aus den beiden Hundertern \(24\) und \(25\) ergibt sich \(36{,}28\); aus den drei \(61,62\) und \(63\) erhält man \(58{,}50\); aus den fünf \(91\ldots95\) ergibt sich \(71{,}56\); schließlich aus den fünf \(96\ldots100\) der Wert \(73{,}54\). Diese Beispiele zeigen, dass die durchschnittliche Anzahl der Klassen zwar langsamer wächst als die Determinanten, aber viel schneller als die durchschnittliche Anzahl der Geschlechter. Bei einiger Aufmerksamkeit erkennt man, dass sie ziemlich genau im Verhältnis der Quadratwurzeln der mittleren Determinanten wächst. Tatsächlich haben wir durch eine theoretische Untersuchung gefunden, dass die durchschnittliche Anzahl der Klassen um die Determinante \(-D\) näherungsweise durch
\[\gamma\surd D-\delta\]
ausgedrückt wird. Dabei ist
\[\gamma=0{,}7467183115=\frac{2\pi}{7e}\]
wobei \(e\) die Summe der Reihe
\[1+\frac18+\frac1{27}+\frac1{64}+\frac1{125}\;\text{etc.}\]
bezeichnet.

Außerdem ist
\[\delta=0{,}2026423673=\frac2{\pi\pi}\]
Die nach dieser Formel berechneten Durchschnittswerte unterscheiden sich wenig von denen, die wir oben aus der Klassifikationstafel entnommen haben. Mit dieser Formel lässt sich auch die Summe der Anzahlen sämtlicher eigentlich primitiven Klassen, bei negativer Determinante der positiven, für
\SourceLeaf{371}
aufeinanderfolgende Determinanten \(-D,-(D+1),-(D+2)\ldots-(D+m-1)\) sehr nahe angeben, wie verschieden die äußersten auch sein mögen. Man addiert die ihnen nach der Formel zugehörigen Durchschnittswerte und erhält
\[=\gamma(\surd D+\surd(D+1)+\text{etc.}+\surd(D+m-1))-\delta m\]
oder sehr nahe
\[=\frac23\gamma((D+m-\frac12)^{\frac32}-(D-\frac12)^{\frac32})-\delta m\]
So berechnet sich zum Beispiel diese Summe für die hundert Determinanten \(-1\ldots-100\) nach der Formel \(=481{,}1\), während sie tatsächlich \(477\) beträgt. Die tausend Determinanten \(-1\ldots-1000\) liefern nach der Tafel \(15533\) Klassen, die Formel gibt \(15551{,}4\). Der zweite Tausender enthält nach der Tafel \(28595\) Klassen, die Formel gibt \(28585{,}7\); ebenso liefert der dritte Tausender tatsächlich \(37092\) Klassen, die Formel \(37074{,}3\); der zehnte Tausender gibt nach der Tafel \(72549\), nach der Formel \(72572\).
